% ==============================================================================
% preamble.tex -- packages, layout, environments and notation for guide.tex
% ==============================================================================

% ---------------------------------------------------------------- fonts & math
\usepackage[T1]{fontenc}
\usepackage[utf8]{inputenc}
\usepackage{newpxtext}
\usepackage{amsmath}
\usepackage{mathtools}
\usepackage{amsthm}
\usepackage[varbb]{newpxmath}
\usepackage[scaled=0.92,varqu,varl]{inconsolata}
% Inconsolata has no italic: typeset code upright inside italic text (theorems, captions)
\AtBeginDocument{%
  \fontencoding{T1}\fontfamily{zi4}\selectfont\normalfont % load T1/zi4 before extending it
  \DeclareFontShape{T1}{zi4}{m}{it}{<->ssub*zi4/m/n}{}%
  \DeclareFontShape{T1}{zi4}{b}{it}{<->ssub*zi4/b/n}{}}
\usepackage{bm}
\usepackage{microtype}
\pdfmapline{=txexs txexs <txexs.pfb} % Type 1 version of the newtx extension symbols (missing from some MiKTeX maps)

% ---------------------------------------------------------------------- layout
\usepackage[a4paper, top=26mm, bottom=28mm, inner=28mm, outer=28mm,
            headheight=15pt, marginparwidth=20mm]{geometry}
\usepackage{setspace}
\setstretch{1.08}
\setlength{\parindent}{0pt}
\setlength{\parskip}{0.55em plus 0.2em minus 0.1em}
\usepackage{emptypage}

% ---------------------------------------------------------- colors & graphics
\usepackage[dvipsnames,table]{xcolor}
\definecolor{primary}{RGB}{16, 62, 120}
\definecolor{primarylight}{RGB}{232, 240, 250}
\definecolor{conv}{RGB}{170, 30, 30}
\definecolor{convlight}{RGB}{253, 240, 240}
\definecolor{warn}{RGB}{200, 110, 0}
\definecolor{warnlight}{RGB}{255, 246, 230}
\definecolor{impl}{RGB}{70, 70, 70}
\definecolor{impllight}{RGB}{245, 245, 245}
\definecolor{proj}{RGB}{20, 110, 70}
\definecolor{projlight}{RGB}{234, 247, 239}
\definecolor{deriv}{RGB}{100, 50, 140}
\definecolor{derivlight}{RGB}{245, 238, 252}
\definecolor{plotblue}{HTML}{0072B2}
\definecolor{plotorange}{HTML}{D55E00}
\definecolor{plotgreen}{HTML}{009E73}
\definecolor{plotpurple}{HTML}{CC79A7}
\definecolor{plotyellow}{HTML}{E69F00}

\usepackage{graphicx}
\graphicspath{{figures/}}
\usepackage{tikz}
\usetikzlibrary{arrows.meta, positioning, calc, decorations.pathreplacing, decorations.pathmorphing,
                shapes.geometric, shapes.misc, fit, backgrounds, patterns, 3d,
                matrix, decorations.markings, shadows}
\tikzset{
  >={Stealth[length=2.2mm]},
  frameaxis/.style={->, line width=1pt},
  block/.style={draw=primary, fill=primarylight, rounded corners=2pt, align=center,
                minimum height=9mm, inner sep=4pt, font=\small},
  note/.style={font=\footnotesize, align=center},
}

% ------------------------------------------------------------- tables & lists
\usepackage{booktabs}
\usepackage{tabularx}
\renewcommand{\tabularxcolumn}[1]{>{\raggedright\arraybackslash}p{#1}} % ragged-right X columns
\usepackage{array}
\usepackage{multirow}
\usepackage{makecell}
\usepackage[shortlabels]{enumitem}
\setlist{itemsep=0.25em, topsep=0.35em, parsep=0.15em}
\usepackage{caption}
\usepackage{subcaption}
\captionsetup{font=small, labelfont={bf,color=primary}, format=hang, width=0.94\linewidth}
\usepackage{float}
\usepackage{placeins}
\usepackage{siunitx}
\sisetup{per-mode=symbol, range-phrase=--, group-digits=integer}
% newpxtext has no text-mode micro sign (\textmu would fall back to a bitmap Computer Modern glyph); use newpxmath's upright mu
\DeclareSIPrefix{\micro}{\ensuremath{\upmu}}{-6}

% -------------------------------------------------------------------- code
\usepackage{listings}
\lstdefinestyle{cpp}{
  language=C++,
  basicstyle=\ttfamily\small,
  keywordstyle=\color{primary}\bfseries,
  commentstyle=\color{proj}\itshape,
  stringstyle=\color{conv},
  backgroundcolor=\color{impllight},
  frame=single, rulecolor=\color{black!15}, framesep=4pt,
  numbers=none, breaklines=true, showstringspaces=false,
  columns=fullflexible, keepspaces=true,
  morekeywords={constexpr, noexcept, nodiscard, consteval, concept, requires, nullptr}
}
\lstdefinestyle{py}{
  language=Python,
  basicstyle=\ttfamily\footnotesize,
  keywordstyle=\color{primary}\bfseries,
  commentstyle=\color{proj}\itshape,
  stringstyle=\color{conv},
  backgroundcolor=\color{impllight},
  frame=single, rulecolor=\color{black!15}, framesep=4pt,
  numbers=none, breaklines=true, showstringspaces=false,
  columns=fullflexible, keepspaces=true
}
\lstset{style=cpp}

% ---------------------------------------------------------------- algorithms
\usepackage{algorithm}
\usepackage[noend]{algpseudocode}
\algrenewcommand\algorithmicrequire{\textbf{Input:}}
\algrenewcommand\algorithmicensure{\textbf{Output:}}
\algnewcommand{\LineComment}[1]{\State \textcolor{proj}{\(\triangleright\) \textit{#1}}}

% ------------------------------------------------------- headings & headers
\usepackage{titlesec}
% Headings are set ragged right and never hyphenated.
\titleformat{\chapter}[display]
  {\normalfont\color{primary}\raggedright\hyphenpenalty=10000}
  {\large\scshape\chaptertitlename\ \thechapter}
  {6pt}
  {\Huge\bfseries}
  [\vspace{3pt}{\color{primary}\titlerule[1.2pt]}]
\titlespacing*{\chapter}{0pt}{-10pt}{26pt}
\titleformat{\section}{\normalfont\Large\bfseries\color{primary}\raggedright\hyphenpenalty=10000}{\thesection}{0.8em}{}
\titleformat{\subsection}{\normalfont\large\bfseries\color{primary!85!black}\raggedright\hyphenpenalty=10000}{\thesubsection}{0.8em}{}
\titleformat{\subsubsection}{\normalfont\normalsize\bfseries}{\thesubsubsection}{0.8em}{}
\titleformat{\part}[display]
  {\centering\normalfont\color{primary}\hyphenpenalty=10000}
  {\LARGE\scshape Part \thepart}
  {14pt}
  {\Huge\bfseries}
\setcounter{secnumdepth}{2}
\setcounter{tocdepth}{1}
\makeatletter % room for two-digit section numbers such as 27.10 in the table of contents
\renewcommand*\l@section{\@dottedtocline{1}{1.5em}{2.9em}}
\makeatother

\usepackage{fancyhdr}
\pagestyle{fancy}
\fancyhf{}
\makeatletter
\renewcommand{\chaptermark}[1]{\markboth{\if@mainmatter\thechapter.\ \fi #1}{}}
\makeatother
\renewcommand{\sectionmark}[1]{\markright{\thesection\ \ #1}}
\fancyhead[L]{\small\color{primary!80!black}\scshape\nouppercase{\leftmark}}
\fancyhead[R]{\small\color{primary!80!black}\nouppercase{\rightmark}}
\fancyfoot[C]{\small\thepage}
\renewcommand{\headrulewidth}{0.3pt}
\fancypagestyle{plain}{\fancyhf{}\fancyfoot[C]{\small\thepage}\renewcommand{\headrulewidth}{0pt}}

% --------------------------------------------------------------------- index
\usepackage[noautomatic]{imakeidx}
\makeindex[intoc, columns=2, title=Index]
\indexsetup{othercode=\small}
% \term{x}: emphasize a term where it is defined and add it to the index.
% \termi{x}{key}: same, with a different index key (e.g. "Kalman filter!gain").
\newcommand{\term}[1]{\emph{#1}\index{#1}}
\newcommand{\termi}[2]{\emph{#1}\index{#2}}

% ------------------------------------------------------------------- version
\newcommand{\guideversion}{1.0}
\newcommand{\guidedate}{October 2026}

% ---------------------------------------------------------------- hyperlinks
\usepackage[pdftex, colorlinks=true, linkcolor=primary, citecolor=proj,
            urlcolor=conv, bookmarksnumbered=true, bookmarksopen=true,
            pdftitle={From Zero to MSCKF: Lecture Notes on Visual-Inertial Odometry (version 1.0)},
            pdfsubject={Visual-inertial odometry, Lie theory, ES-EKF, MSCKF},
            pdfkeywords={version 1.0}]{hyperref}
\usepackage[nameinlink, capitalize, noabbrev]{cleveref}
\usepackage{bookmark}
\makeatletter % unique hyperref anchors for algorithm lines across algorithms
\providecommand{\theHALG@line}{}
\renewcommand{\theHALG@line}{\thealgorithm.\arabic{ALG@line}}
\makeatother

% ------------------------------------------------------ theorem environments
\newtheoremstyle{notesplain}{6pt}{6pt}{\itshape}{}{\bfseries\color{primary}}{.}{0.5em}{}
\newtheoremstyle{notesdef}{6pt}{6pt}{\normalfont}{}{\bfseries\color{primary}}{.}{0.5em}{}
\newtheoremstyle{notesrem}{6pt}{6pt}{\normalfont}{}{\bfseries\itshape\color{primary!80!black}}{.}{0.5em}{}
\theoremstyle{notesplain}
\newtheorem{theorem}{Theorem}[chapter]
\newtheorem{proposition}[theorem]{Proposition}
\newtheorem{lemma}[theorem]{Lemma}
\newtheorem{corollary}[theorem]{Corollary}
\theoremstyle{notesdef}
\newtheorem{definition}[theorem]{Definition}
\newtheorem{example}[theorem]{Example}
\theoremstyle{notesrem}
\newtheorem{remark}[theorem]{Remark}
\crefname{theorem}{Theorem}{Theorems}
\crefname{proposition}{Proposition}{Propositions}
\crefname{lemma}{Lemma}{Lemmas}
\crefname{corollary}{Corollary}{Corollaries}
\crefname{definition}{Definition}{Definitions}
\crefname{example}{Example}{Examples}
\crefname{remark}{Remark}{Remarks}
\renewcommand{\qedsymbol}{\color{primary}$\blacksquare$}

\usepackage[most]{tcolorbox}
\tcbuselibrary{breakable, skins, theorems}
\tcolorboxenvironment{theorem}{blanker, breakable, borderline west={2pt}{0pt}{primary}, left=9pt, before skip=8pt, after skip=8pt}
\tcolorboxenvironment{proposition}{blanker, breakable, borderline west={2pt}{0pt}{primary}, left=9pt, before skip=8pt, after skip=8pt}
\tcolorboxenvironment{lemma}{blanker, breakable, borderline west={2pt}{0pt}{primary!60}, left=9pt, before skip=8pt, after skip=8pt}
\tcolorboxenvironment{corollary}{blanker, breakable, borderline west={2pt}{0pt}{primary!60}, left=9pt, before skip=8pt, after skip=8pt}
\tcolorboxenvironment{definition}{blanker, breakable, colback=primarylight, left=9pt, right=6pt, top=4pt, bottom=4pt, borderline west={2pt}{0pt}{primary}, before skip=8pt, after skip=8pt}
\tcolorboxenvironment{example}{blanker, breakable, borderline west={1pt}{0pt}{proj}, left=9pt, before skip=8pt, after skip=8pt}

% ------------------------------------------------------------- callout boxes
\newtcolorbox{keyidea}[1][]{enhanced, breakable, colback=primarylight, colframe=primary,
  boxrule=0.7pt, arc=1.5mm, left=8pt, right=8pt, top=5pt, bottom=5pt,
  fonttitle=\bfseries\small, title={Key idea\ifx\relax#1\relax\else: #1\fi},
  before skip=10pt, after skip=10pt}
\newtcolorbox{convention}[1][]{enhanced, breakable, colback=convlight, colframe=conv,
  boxrule=0.9pt, arc=1.5mm, left=8pt, right=8pt, top=5pt, bottom=5pt,
  fonttitle=\bfseries\small, title={Frozen convention\ifx\relax#1\relax\else: #1\fi},
  before skip=10pt, after skip=10pt}
\newtcolorbox{pitfall}[1][]{enhanced, breakable, colback=warnlight, colframe=warn,
  boxrule=0.7pt, arc=1.5mm, left=8pt, right=8pt, top=5pt, bottom=5pt,
  fonttitle=\bfseries\small, title={Pitfall\ifx\relax#1\relax\else: #1\fi},
  before skip=10pt, after skip=10pt}
\newtcolorbox{implnote}[1][]{enhanced, breakable, colback=impllight, colframe=impl,
  boxrule=0.6pt, arc=1.5mm, left=8pt, right=8pt, top=5pt, bottom=5pt,
  fonttitle=\bfseries\small, title={Implementation note\ifx\relax#1\relax\else: #1\fi},
  before skip=10pt, after skip=10pt}
\newtcolorbox{projectlink}[1][]{enhanced, breakable, colback=projlight, colframe=proj,
  boxrule=0.6pt, arc=1.5mm, left=8pt, right=8pt, top=4pt, bottom=4pt,
  fonttitle=\bfseries\small, title={Project link\ifx\relax#1\relax\else: #1\fi},
  before skip=10pt, after skip=10pt}
\newtcolorbox{objectives}{enhanced, breakable, colback=white, colframe=primary!60,
  boxrule=0.5pt, arc=0mm, left=8pt, right=8pt, top=4pt, bottom=4pt,
  fonttitle=\bfseries\small\color{primary}, coltitle=primary, colbacktitle=primarylight,
  title={In this chapter}, before skip=4pt, after skip=12pt}
\newtcolorbox{takeaways}{enhanced, breakable, colback=primarylight!60, colframe=primary,
  boxrule=0.8pt, arc=1.5mm, left=8pt, right=8pt, top=5pt, bottom=5pt,
  fonttitle=\bfseries, title={Key takeaways}, before skip=14pt, after skip=10pt}

\newtcolorbox{aside}[1][]{enhanced, breakable, colback=white, colframe=primary!45,
  boxrule=0.5pt, arc=1.5mm, left=8pt, right=8pt, top=5pt, bottom=5pt, borderline west={2pt}{0pt}{primary!45},
  fonttitle=\bfseries\small, coltitle=primary!80!black, colbacktitle=white, attach boxed title to top left={xshift=4mm, yshift=-2mm},
  boxed title style={colframe=white, colback=white}, title={Going deeper\ifx\relax#1\relax\else: #1\fi},
  before skip=12pt, after skip=10pt}
\newtcolorbox{worked}[1][]{enhanced, breakable, colback=projlight!45, colframe=proj!70!black,
  boxrule=0.6pt, arc=1.5mm, left=8pt, right=8pt, top=5pt, bottom=5pt,
  fonttitle=\bfseries\small, title={Worked example\ifx\relax#1\relax\else: #1\fi},
  before skip=10pt, after skip=10pt}

% Derivation exercises for the project Jacobians (AGENTS.md Guardrail 4).
% Full solutions live in the appendix, which is compiled only when \ifsolutions.
\newif\ifsolutions
\ifdefined\NoSolutions \solutionsfalse \else \solutionstrue \fi
\newtcolorbox{derivex}[2]{enhanced, breakable, colback=derivlight, colframe=deriv,
  boxrule=0.9pt, arc=1.5mm, left=8pt, right=8pt, top=5pt, bottom=5pt,
  fonttitle=\bfseries\small, title={Derivation exercise (project Jacobian): #1},
  before skip=10pt, after skip=10pt,
  after upper={\par\smallskip\hfill\footnotesize\itshape
    \ifsolutions Worked solution: \cref{#2} (page~\pageref{#2}).\else
    The worked solution is omitted from this edition; build without \texttt{NoSolutions} to include it.\fi}}

% Exercise lists at chapter ends, and their worked solutions on the following pages.
% Each exercise number links to its solution and each solution number links back.
\newcounter{exercisectr}[chapter]
\newenvironment{exercises}{%
  \section*{Exercises}\addcontentsline{toc}{section}{Exercises}\markright{Exercises}%
  \begin{enumerate}[label={\protect\hypertarget{ex.\thechapter.\arabic*}{}%
      \protect\hyperlink{sol.\thechapter.\arabic*}{\color{primary}\bfseries\thechapter.\arabic*}},
    leftmargin=2.6em, itemsep=0.55em]}%
  {\end{enumerate}}
\newenvironment{solutions}{%
  \clearpage
  \section*{Solutions to the exercises}\addcontentsline{toc}{section}{Solutions to the exercises}%
  \markright{Solutions to the exercises}%
  \begin{enumerate}[label={\protect\hypertarget{sol.\thechapter.\arabic*}{}%
      \protect\hyperlink{ex.\thechapter.\arabic*}{\color{primary}\bfseries\thechapter.\arabic*}},
    leftmargin=2.6em, itemsep=0.9em]}%
  {\end{enumerate}}
% A solution that would reveal part of a project Jacobian (AGENTS.md Guardrail 4) is printed only
% in the solutions edition: \item \ifsolutions <solution> \else \withheldsolution \fi
\newcommand{\withheldsolution}{\textit{This solution contains part of a project Jacobian (AGENTS.md Guardrail 4)
  and is omitted from this edition. The finite-difference test of the corresponding derivation exercise tells you
  whether your own answer is right.}}

% ------------------------------------------------------------------ notation
\newcommand{\R}{\mathbb{R}}
\newcommand{\N}{\mathcal{N}}
\newcommand{\E}{\mathbb{E}}
\newcommand{\T}{^{\mathsf{T}}}
\newcommand{\mT}{^{-\mathsf{T}}}
\newcommand{\inv}{^{-1}}
\newcommand{\sx}[1]{[#1]_{\times}}
\newcommand{\qprod}{\otimes}
\newcommand{\Dt}{{\Delta t}} % braced: safe in sub/superscripts, \sqrt and \frac
\newcommand{\dd}{\mathrm{d}}
\newcommand{\eps}{\varepsilon}
\newcommand{\SO}{\mathrm{SO}(3)}
\newcommand{\so}{\mathfrak{so}(3)}
\newcommand{\SE}{\mathrm{SE}(3)}
\newcommand{\se}{\mathfrak{se}(3)}
\newcommand{\Jr}{J_r}
\newcommand{\Jl}{J_l}
\newcommand{\Jrinv}{J_r^{-1}}
\newcommand{\Jlinv}{J_l^{-1}}
\newcommand{\dth}{{\delta\theta}}
\newcommand{\dx}{{\delta x}}
\newcommand{\bplus}{\boxplus}
\newcommand{\bminus}{\boxminus}
\newcommand{\code}[1]{{\upshape\ttfamily #1}}
\newcommand{\hatop}{^{\wedge}}
\newcommand{\veeop}{^{\vee}}
\newcommand{\mach}{\varepsilon_{\mathrm{mach}}}
\DeclareMathOperator{\Exp}{Exp}
\DeclareMathOperator{\Log}{Log}
\DeclareMathOperator{\tr}{tr}
\DeclareMathOperator{\diag}{diag}
\DeclareMathOperator{\blkdiag}{blkdiag}
\DeclareMathOperator{\rank}{rank}
\DeclareMathOperator{\Cov}{Cov}
\DeclareMathOperator{\Var}{Var}
\DeclareMathOperator{\atantwo}{atan2}
\DeclareMathOperator{\vect}{vec}
\DeclareMathOperator{\Ad}{Ad}
\DeclareMathOperator{\ad}{ad}
\DeclareMathOperator{\range}{range}
\DeclareMathOperator{\nullsp}{null}
\DeclareMathOperator*{\argmin}{arg\,min}
\DeclareMathOperator*{\argmax}{arg\,max}
\DeclareMathOperator{\sinc}{sinc}

\newcommand{\figref}[1]{\cref{#1}}
